% Extracto íntegro por residencia del propietario científico definitivo de masas.
\section{Equivalencia entre la sección persistente y su representación finita}

La definición genealógica de partícula y su representación operatorial son
dos presentaciones del mismo objeto. Sea
\[
 \mathfrak r_K=(x_K,\gamma_K,B_K,s_K,n_K,U_K)
\]
un representante finito a profundidad \(K\), con estado, ruta, frontera,
orientación, contador y representación, y sean
\(r_{L,K}:\mathfrak R_L\to\mathfrak R_K\) los mapas de refinamiento. Una
sección persistente es una familia
\[
 X=(\mathfrak r_K)_{K\ge K_0},
 \qquad r_{L,K}(\mathfrak r_L)=\mathfrak r_K
 \quad(L\ge K\ge K_0).
\]
Dos familias son equivalentes si coinciden en una cola cofinal.

\begin{theorem}[Equivalencia de la sección persistente y el límite inverso]
La categoría de secciones persistentes, definidas como familias compatibles,
es canónicamente equivalente a
\[
 \varprojlim_K\mathfrak R_K/\!\sim_{\rm cola}.
\]
La proyección a una profundidad finita es una evaluación, no un funtor inverso.
Un representante \(\mathfrak r_K\) determina una única sección persistente si
y sólo si posee una única prolongación compatible módulo equivalencia de cola;
esta unicidad se obtiene, por ejemplo, cuando los mapas de refinamiento son
biyectivos sobre la cola superviviente a partir de \(K\).
\end{theorem}

\begin{proof}
Por definición, una sección persistente es precisamente un punto del límite
inverso; el cociente identifica las familias que coinciden en una cola cofinal.
La evaluación es la proyección canónica a una componente. Su fibra contiene
todas las prolongaciones compatibles del representante evaluado, de modo que
es unitaria exactamente bajo la condición de prolongación única indicada.
\end{proof}

Una partícula HMT--MD es, por tanto, una clase de secciones persistentes sin
obstrucción terminal finita, provista de representación, orientación, carga,
espín y carácter espectral compatibles con el refinamiento. La ruta es su
historia observable, no su identidad reducida a una celda. La antipartícula
es la sección conjugada por la involución dodecafásica y la representación
conjugada. Una sección finita obstruida es virtual; un producto cocíclico de
secciones es compuesto; una representación trenzada persistente es anyónica.

\section{Teorema de cierre tipado de la realización másica}

\begin{theorem}[Cierre por codominio físico]
Supónganse construidos el atlas \(81/56/13/324\), el sector nulo, la escala
acción--reloj--masa, los lectores \(K_D,K_\Omega\), el descriptor físico y la
representación interna de cada estado. Entonces todo descriptor admisible y
realizable, es decir, todo \(X\) con
\(d_{\mathcal P}(X)\in\operatorname{im}d_{\mathcal R}\),
determina, sin consultar su valor externo:
\begin{enumerate}
\item una fibra, órbita o multisección \(\mathscr S(X)\);
\item el par positivo bidireccional y su PVM conjunta
      \((\widehat M_X^{\beta,D},\widehat M_X^{\beta,\Omega},
      E_X^{D,\Omega})\);
\item una descomposición en bloques irreducibles;
\item una medida espectral para cada acoplamiento físico;
\item un escalar en un sector comprimido de rango uno o un multiplete en un
      sector reducible; en un sector abierto, un operador efectivo y, sólo si
      existe continuación meromorfa con un cero aislado admisible, un polo
      complejo;
\item un objeto categórico cuando la realización es topológica y no posee aún
      una brecha Hamiltoniana;
\item un criterio de inexistencia cuando el tipo solicitado contradice
      positividad, calibre o asintoticidad.
\end{enumerate}
Ninguno de estos resultados requiere seleccionar una ruta mediante la masa
que después será contrastada.
\end{theorem}

\begin{proof}
El producto fibrado construye la fibra de forma prospectiva. El carácter y los
dos lectores construyen el par positivo y su PVM conjunta. Cuando el tipo
físico declara un funcional escalar, el cálculo funcional produce el operador
correspondiente; la esperanza condicional lo proyecta al conmutante de la
simetría. La descomposición isótípica produce escalares o bloques por el lema
de Schur. Los proyectores de preparación producen medidas espectrales. La
reducción de Feshbach produce el operador efectivo; la continuación meromorfa
y un cero aislado admisible producen un polo, si existe. Los sectores
topológicos conservan sus datos de fusión y trenzado hasta que exista una
realización Hamiltoniana. La positividad y las restricciones de calibre
excluyen los tipos incompatibles. Cada paso depende sólo de datos internos o
del acoplamiento declarado, nunca de la columna externa.
\end{proof}

\section{Condiciones de consistencia e incompatibilidad de las realizaciones espectrales, fermiónicas y compuestas}

La construcción queda refutada en su dominio si ocurre cualquiera de los
hechos siguientes:

\begin{enumerate}
\item al modificar los valores externos cambia la fibra \(\mathscr S(X)\);
\item la asignación deja de ser equivariante bajo inversión celular,
      conjugación o calibre;
\item se selecciona una celda puntual en una multisección no central sin un
      acoplamiento que rompa la simetría;
\item la síntesis de \(K_D\) y \(K_\Omega\) privilegia una orientación o no
      recupera la diagonal cuando ambos lectores coinciden;
\item una compresión con varios autovalores se reemplaza por una media sin
      proyector físico;
\item el cociente de color produce masas diferentes para rojo, verde y azul;
\item se comparan masas quark sin declarar el mismo esquema y la misma escala;
\item una fase de rango uno se presenta como mezcla PMNS después de fallar el
      contraste angular de rango completo;
\item la neutralidad se usa por sí sola para afirmar el carácter Majorana;
\item la masa de un compuesto se obtiene sumando masas de constituyentes sin
      registro de enlace;
\item la anchura se inserta como coordenada de la firma en vez de extraerse de
      un modo abierto o un polo;
\item un autovalor negativo del Hessiano se publica como partícula taquiónica;
\item la gravitación se hace depender de un gravitón elemental introducido
      antes de la conexión, la curvatura y la torsión.
\end{enumerate}

\section{Contraste metrológico y archivo exhaustivo de realizaciones}

La comparación externa se abre únicamente después de congelar descriptor,
fibra, operador, proyector, esquema y escala. Cada fila debe distinguir masa
estable, masa de polo, estimador de resonancia y anchura. El atlas de \(324\)
fichas y los inventarios de \(471\) masas y \(384\) anchuras constituyen la
exposición exhaustiva subsiguiente. La construcción conceptual y cada ficha
se remiten recíprocamente mediante sus identificadores de ruta. Las fichas
conservan todos sus campos y se componen como unidades coherentes; no se
sustituyen por tablas históricas abreviadas.

La comparación no consiste en ampliar una lista de nombres, sino en aplicar,
para cada tipo físico, el mapa completo que conduce desde la genealogía de ruta
hasta su codominio observable.
